\chapter{Summary -- Kết luận}\label{ch:summary}
Luận văn xem trajectory privacy như một bài toán bảo vệ transcript dịch vụ theo thời gian. Một
tọa độ giả ở từng thời điểm chưa đủ: quyền quan sát cả chuỗi, quyết định cập nhật và tương quan
giữa các đầu ra đều cần được đặc tả. Đồng thời, utility được đánh giá bằng đáp án POI người dùng
thực nhận thay vì chỉ độ lệch giữa GPS và tọa độ phát.

\section{Trả lời các câu hỏi nghiên cứu}
\textbf{RQ1 -- Bảo đảm privacy toàn transcript.} GeoI-BQ v1 nối REM/noisy reuse với một prospective
cap cho private GPS và một guard riêng cho lựa chọn đoạn Q. Trên miền/support chung và các giả thiết
kernel, clock, ledger, provider đã nêu, chuỗi Q có tỷ số likelihood không quá $R_n\le e^{\Gamma_Q}$;
với $\Gamma_Q=1$, cận toàn thời lượng nhỏ hơn 2,437372. Nhánh Geo-I lý tưởng cho cận theo metric
trajectory; khi cả hai nhánh hợp lệ có thể lấy minimum của chúng. Đây là giới hạn mức transcript
làm thay đổi khả năng phân biệt/suy luận, không là chứng minh attacker không thể biết vị trí.

\textbf{RQ2 -- Utility và các điều kiện cần.} Centered kernel có gain không âm trên planning score,
còn actual utility gain phụ thuộc covariance/proxy calibration. Union/cache và local ranking giữ
được các reference POI đã truy hồi tại cùng GPS/request khi metadata hợp lệ. Phân tích toàn history
cho cận utility so matched public base bằng conditional expectation, không dùng IID. Floor chuyển
sang GPS thật vẫn cần bound sai số và provider completeness. Những điều kiện chưa chứng nhận không
được đặt bằng không trong kết luận.

\textbf{RQ3 -- Trade-off thực nghiệm.} Trong controlled comparison với ba adaptation đã audit từ
paper và Raw/REM-Z controls, GeoI-BQ đạt session Recall@5 99,595\%, CI95\% [99,224;99,903]. Các
coordinate tasks S1--S3/S9--S10 có Hit100 0--0,188\% dưới bank đã chọn, thấp hơn các adaptation trong
trial, đổi lại dùng 155 point queries mỗi phiên và utility thấp hơn khoảng 0,4 điểm phần trăm so
ba adaptation. Các con số không chứng minh mọi original paper implementation hoặc mọi attacker
đều thua, và không là generalization ngoài cohort.

\section{Đóng góp thực sự của bản hiện tại}
Thứ nhất, luận văn hệ thống hóa targets, output contracts và một protocol đánh giá bổ sung giữa
các cơ chế khác dạng. MAE, Recall và bootstrap là kiến thức kế thừa; đóng góp là cách ghép target,
window, aggregation, cost và failure accounting thành một phép đối chiếu kiểm tra được.

Thứ hai, phương pháp tách protected reference khỏi public query action, quản lý noisy reuse theo
phiên và chọn chung đoạn Q bằng kernel bị chặn có integer representation. Các chứng minh liên kết
từ thành phần đến transcript và nêu rõ điều kiện của suy luận Bayes cũng như utility. Không quy
đóng góp thành việc tạo thêm nhiễu hoặc phát minh lại Geo-I.

Thứ ba, artifact chứa main comparison, các tác vụ S4--S8 có scope riêng, bảng CIs/selected attackers,
source pins và verifier. Quá trình sửa attacker S8 được giữ để kết quả có thể bị kiểm tra và phản
biện. Đây là bằng chứng bổ sung cho thiết kế, không thay cho các giả thiết toán học.

\section{Giới hạn và hướng phát triển}
Ưu tiên tiếp theo là chứng nhận sampler upstream/restore toàn state và nguồn randomness; sau đó
đánh giá trên cohort/bản đồ mới và catalogue lớn/động với latency, bytes metadata đầy đủ và energy
thực đo. Matched public base đang có utility tốt hơn trên cohort hiện tại, nên cần xác định điều
kiện nào private adaptation thực sự đáng chi phí thay vì mặc định $\Gamma=1$ tối ưu.

Identity người/xe cần nhãn vật lý và task riêng; route-family association chưa đủ. S5 cần cadence
và positive control phù hợp; S6 cần lựa chọn/historical context đa dạng hơn. S7 cần đo route-correlated
intent ngoài payload, S8 cần dữ liệu tương quan thực và phân tích joint channels/multiple sessions.
Các hướng đó mở rộng phạm vi đã có, không được trình bày như kết quả hoàn thành trong bản này.

Bản hiện tại cung cấp một cơ sở thống nhất để review: kiến trúc/đặc tả trước, chứng minh có điều kiện
tiếp theo, rồi benchmark cùng protocol và kết luận tương xứng. Từ đó có thể chỉnh sửa diễn giải,
củng cố chứng nhận triển khai và xác nhận ngoài miền mà vẫn giữ nền tảng Geo-I của phương pháp.
